\section{Prompts}
\label{app:prompts}
The system prompt of Developer 1 on an instance of the core regime (abstract rendering), with the \neutral{} guidance. Scores and the outside option are the agent's own; the partner's prompt mirrors it. Under a skin, ``Developer 1'' becomes the skin's party (``Team lead 1''), decisions and options carry names (``Hiring: A (two senior engineers) = 48''), and ``design'' becomes the skin's artifact; a grounded instance adds the context paragraph and each side's ticket (``You are working on \emph{Add batch slug checks and user lookup APIs}; the other developer on \emph{Short-circuit scalar namespace checks}''). Each move is a user turn that quotes the partner's message and proposal with its value in the agent's own points, followed by ``Move $n$ of $T$. Your move.''

{\fontsize{8}{9.5}\selectfont\begin{verbatim}
You are Developer 1, one of two developers who must agree on a shared design with 4 decisions.
For each decision exactly one option will be chosen. Your private scores for each option are
below; a design is worth the sum of your scores for its chosen options, at most 100 points.
The other developer has their own private scores, which you do not see.
If you both accept one complete design, each of you gets your own value of it. If no design is
accepted within 10 moves in total, or either of you walks away, you each get your outside
option: yours is 30 points. Your goal is to get as many points as you can.

Your scores:
Decision 1: Option A = 48, Option B = 24, Option C = 0
Decision 2: Option A = 0, Option B = 10.5, Option C = 21
Decision 3: Option A = 8, Option B = 0, Option C = 16
Decision 4: Option A = 7.5, Option B = 0, Option C = 15

Moves alternate. On your move, reply with one JSON object and nothing else, in one of these forms:
{"message": "<what you say to the other developer>", "action": "propose",
 "proposal": {"Decision 1": "Option A", "Decision 2": "Option A", "Decision 3": "Option A",
              "Decision 4": "Option A"}}
{"message": "<...>", "action": "accept"}   (accepts the other developer's latest proposal as stated)
{"message": "<...>", "action": "walk"}   (ends the negotiation; both get their outside options)
A proposal names one option for every decision. You may say anything in the message, including
your scores, or nothing. You do not have to accept a proposal you consider poor: you can
counter-propose, and the other developer can accept your proposal. There are 10 moves in total.

[Guidance] You and the other developer want different things from this design: each of you
wants the options that score well for your own side. Handle this in whatever way you think best.
\end{verbatim}}

